\documentclass[12pt]{article}
\usepackage[spanish]{babel}
\usepackage[utf8]{inputenc}
\usepackage[T1]{fontenc}
\usepackage{amsmath, amssymb}
\usepackage{geometry}
\usepackage{enumitem}
\usepackage{needspace}
\usepackage{tikz}
\usetikzlibrary{babel}

\geometry{letterpaper, margin=2cm}
\setlength{\parindent}{0pt}
\setlength{\parskip}{6pt}

\newcommand{\puntografico}{0.35}
\newcommand{\opcionalgebra}[2]{\item[#1)] $#2$}
\newcommand{\graftrasbase}[7]{%
\begin{tikzpicture}[scale=0.16]
  \path[use as bounding box] (-11,-9) rectangle (31,24);
  \draw[<->, thick] (-8,0) -- (28,0);
  \draw[<->, thick] (0,-7) -- (0,21);
  \node[anchor=west] at (29,-0.7) {$x$};
  \node[anchor=south east] at (3,22) {$y$};
  \draw[thick] (#4,#5) circle (#3);
  \draw[dashed] (#4,0) -- (#4,#5);
  \draw[dashed] (0,#5) -- (#4,#5);
  \fill (#4,#5) circle (\puntografico) node[anchor=south west, inner sep=2pt] {$A'$};
  #7
  \node[left] at (0,#5) {$#5$};
\end{tikzpicture}}
\newcommand{\graftras}[6]{%
\graftrasbase{#1}{#2}{#3}{#4}{#5}{#6}{\node[anchor=north] at (#4,0) {$#4$};}}
\newcommand{\graftrasetiquetabaja}[6]{%
\graftrasbase{#1}{#2}{#3}{#4}{#5}{#6}{\node[anchor=north, inner sep=1pt] at (#4,#5-0.45) {$#4$};}}

\begin{document}

\begin{center}
  {\Large \textbf{Borrador de items}}\\
  {\large Bloque 1. Geometria - Afirmacion 2}\\[4pt]
  Prueba Nacional Estandarizada de Matematica, Secundaria 2026
\end{center}

\section*{Afirmacion}
Resuelve problemas relacionados con traslaciones de circunferencias.

\section*{Items propuestos}

\begin{enumerate}[label=\textbf{\arabic*.}, itemsep=1.05cm]

\Needspace{0.68\textheight}
\item
En un plano cartesiano, la circunferencia $C$ tiene centro en $A(6,4)$ y radio $5$. A esa circunferencia se le aplica una traslacion de $8$ unidades hacia la derecha y $3$ unidades hacia abajo, obteniendo la circunferencia $C'$.

\textquestiondown Cual de las siguientes representaciones algebraicas corresponde a la circunferencia $C'$?

\begin{enumerate}[label=\Alph*)]
  \opcionalgebra{A}{(x-14)^2+(y-1)^2=25}
  \opcionalgebra{B}{(x+14)^2+(y+1)^2=25}
  \opcionalgebra{C}{(x-14)^2+(y-7)^2=25}
  \opcionalgebra{D}{(x-6)^2+(y-4)^2=25}
\end{enumerate}

\Needspace{0.90\textheight}
\item
En un centro recreativo, una piscina circular esta representada en un plano por la circunferencia $C$, con centro en $A(4,3)$ y radio $4$. Debido a una remodelacion del terreno, la piscina se trasladara $12$ unidades hacia la derecha y $5$ unidades hacia arriba, obteniendo la nueva ubicacion representada por la circunferencia $C'$.

\textquestiondown Cual de las siguientes representaciones graficas corresponde a esa traslacion?

\vspace{2\baselineskip}

\begin{center}
\setlength{\tabcolsep}{0pt}
\renewcommand{\arraystretch}{1}
\begin{tabular}{@{}r@{\hspace{0.35cm}}c@{\hspace{0.55cm}}r@{\hspace{0.35cm}}c@{}}
\raisebox{1.45cm}{\textbf{A)}} & \graftras{4}{3}{4}{12}{8}{4} &
\raisebox{1.45cm}{\textbf{B)}} & \graftrasetiquetabaja{4}{3}{4}{16}{-2}{4} \\[0.75cm]
\raisebox{1.45cm}{\textbf{C)}} & \graftras{4}{3}{4}{16}{8}{4} &
\raisebox{1.45cm}{\textbf{D)}} & \graftras{4}{3}{4}{9}{15}{4}
\end{tabular}
\end{center}

\Needspace{0.72\textheight}
\item
En una plaza, el borde de una zona circular de juego esta representado por la circunferencia
\[
(x-18)^2+(y-10)^2=36,
\]
donde las unidades estan en metros. Por una remodelacion, toda la zona se trasladara $7$ m hacia el oeste y $4$ m hacia el norte.

\textquestiondown Cual sera la representacion algebraica del borde de la zona circular despues de la traslacion?

\begin{enumerate}[label=\Alph*)]
  \opcionalgebra{A}{(x-11)^2+(y-14)^2=36}
  \opcionalgebra{B}{(x-25)^2+(y-14)^2=36}
  \opcionalgebra{C}{(x-11)^2+(y-6)^2=36}
  \opcionalgebra{D}{(x-18)^2+(y-10)^2=36}
\end{enumerate}

\end{enumerate}

\section*{Clave y justificacion breve}

\begin{enumerate}[label=\textbf{\arabic*.}]
  \item \textbf{A.} El centro se traslada de $(6,4)$ a $(14,1)$ y el radio conserva su medida, por eso $r^2=25$.
  \item \textbf{C.} El centro imagen es $A'(16,8)$, pues $(4+12,3+5)=(16,8)$; el radio no cambia.
  \item \textbf{A.} El centro original es $(18,10)$; al trasladarlo $7$ m al oeste y $4$ m al norte se obtiene $(11,14)$, con el mismo radio.
\end{enumerate}

\end{document}
